\documentclass[10pt,a4paper]{amsart}
\usepackage[margin=19mm]{geometry}
\usepackage{amsmath,amssymb,mathtools}
\usepackage[scr=rsfs,cal=euler]{mathalfa}
\usepackage{xfrac}
\usepackage[unicode,hidelinks,bookmarks=false]{hyperref}
\newcommand{\ie}{i.e.\ }
\newcommand{\cf}{cf.\ }
\newcommand{\resp}{respectively\ }
\newcommand{\Subsection}{\subsection}
\providecommand{\nobreakdash}{\nobreak\hbox{-}}
\setlength{\emergencystretch}{2em}
\multlinegap0pt
\usepackage{bm}
\usepackage{bbm}
\usepackage{dsfont}
\usepackage{xspace}
\usepackage{cancel}
\usepackage{mathtools}
\usepackage{amsthm}
\makeatletter
\@ifundefined{thm}{\newtheorem{thm}{Theorem}}{}
\makeatother
\makeatletter
\expandafter\ifx\csname enumerit\endcsname\relax
\newenvironment{enumerit}{\begin{list}{{\hfill\rm(\roman{cnt})\hfill}}{%
\settowidth{\labelwidth}{{\rm(iv)}}\leftmargin=\labelwidth%
\advance\leftmargin by \labelsep\rightmargin=0pt\usecounter{cnt}}}{\end{list}} \makeatletter
\fi
\expandafter\ifx\csname pf\endcsname\relax
\newenvironment{pf}{\proof}{\endproof}
\fi
\makeatother
\title[Multivariate growth series of graph products of groups]{Multivariate growth series of graph products of groups}
\author{Chaithra Pilakkat, Venkatesh Rajendran}
\date{Source: arXiv:2607.23668v1 (CC BY 4.0)}
\begin{document}
\maketitle

\noindent\emph{Source and licence.} Multivariate growth series of graph products of groups. Version arXiv:2607.23668v1 (CC BY 4.0), distributed under the Creative Commons Attribution 4.0 International licence, as recorded in the arXiv licence metadata for this exact version and shown on the abstract page https://arxiv.org/abs/2607.23668v1. Full rights notice: licenses/arxiv-2607.23668-CC-BY-4.0.txt.\par\smallskip

\noindent\textit{Editorial scope.} Counted unit: the Thm whose source label is \texttt{mainthm} in the article (source0.tex lines 354--363, source label \texttt{mainthm}), delivered together with the article's own complete original proof of it (source0.tex lines 366--493, one \texttt{proof} environment of 128 lines). No mathematics is written, repaired or re-derived: statement and proof are the article's own text line for line. Required context retained uncounted: none. Every definition, display and label of those lines is retained unchanged. Omitted from this package and not redistributed: 1--353, 364--365, 494--1417. Nothing inside a retained excerpt is omitted or altered: every retained body is delivered exactly as the article's lines are written, so no omission receipt of any kind accompanies this package. Citations: the retained excerpts carry no citation command at all, so no bibliography entry of the article is transcribed and no reference is redistributed. Reference adaptation: none was needed and none was made; every reference inside the delivered unit resolves inside it, so no unresolved reference or citation is produced. Printed slips kept literal: no printed word, symbol, hypothesis or number of the counted statement or of its proof was altered. Layout and class: the article's own class is \texttt{amsart} with packages including amsthm, amssymb, latexsym, multicol, verbatim, enumerate, color, xy, inputenc, hyperref, amsmath, amscd, soul, tikz; the wrapper is the lane's pinned \texttt{amsart} editorial wrapper at 10pt on A4, which restates the theorem-like environments the retained text needs and adds the article's own macro definitions that the retained text needs (none), with extra wrapper packages (bm, bbm, dsfont, xspace, cancel, mathtools). The delivered numbering is the unit's own. Assets: no retained span contains a figure environment, a graphics-inclusion command, a PDF-page inclusion, a table or a TikZ picture; no image file is copied into this package, the package declares no asset dependency and no image file is redistributed with it. Licence: version arXiv:2607.23668v1 of this article is distributed under the Creative Commons Attribution 4.0 International licence, as recorded in the arXiv licence metadata for this exact version and shown on the abstract page https://arxiv.org/abs/2607.23668v1. A full rights notice is delivered at licenses/arxiv-2607.23668-CC-BY-4.0.txt. Characters: every retained excerpt is pure ASCII, so the delivered engine, XeLaTeX with its default Latin Modern text font, reports no missing character.
\begin{thm}\label{mainthm}
    Let $\mathcal{G}=(I,E)$ be a simple graph, and let $\mathcal{I}(\mathcal{G})$ denotes the set of all stable subsets of $\mathcal{G}$. The multivariate growth series $H_{\Gamma}(\mathbf{x})$ of the graph product $\Gamma =  \Gamma(\mathcal G, \{G_i\}_{i\in I}) $ is given by
\begin{equation}\label{equationthm}
H_{\Gamma}(\mathbf{x}) = \left( \sum_{S \in \mathcal{I}(\mathcal{G})} \prod_{i \in S} \left( \frac{1}{H_i(x_i)} - 1 \right) \right)^{-1}.
\end{equation}
%\textit{Equivalently,}
%\begin{equation*}
%H_{\Gamma}(\mathbf{x}) = \left( \sum_{S \in\mathcal{I}(\mathcal{G})} \prod_{i \in S} \frac{-F_i(x_i)}{1 + F_i(x_i)} \right)^{-1}.
%\end{equation*}
\end{thm} 

\begin{proof}
    Set
\begin{equation*}
D(\mathbf{x}) = \sum_{S \in \mathcal I(\mathcal G)} \prod_{i \in S} \left( \frac{1}{H_i(x_i)} - 1 \right).
\end{equation*}
It suffices to prove
\begin{equation*}
H_{\Gamma}(\mathbf{x})D(\mathbf{x}) = 1.
\end{equation*}
%Expand
%\begin{equation*}
%\frac{1}{H_i(x_i)} - 1 = \frac{-F_i(x_i)}{1 + F_i(x_i)} = \sum_{r \ge 1} (-1)^r F_i(x_i)^r.
%\end{equation*}
Recall that we have $F_i(x_i) = H_i(x_i) - 1$ and $\frac{1}{H_i(x_i)} - 1 = \frac{-F_i(x_i)}{1 + F_i(x_i)} = \sum_{r \ge 1} (-1)^r F_i(x_i)^r$. 
A monomial of $F_i(x_i)^r$ (for $r\ge1$) corresponds to an ordered tuple
\begin{equation*}
u_i = (h_{i,1}, \dots, h_{i,r}), \quad h_{i,j} \in G_i \setminus \{1\}.
\end{equation*}
Thus a term of
\begin{equation*}
H_{\Gamma}(\mathbf{x})D(\mathbf{x})
\end{equation*}
is represented by a pair $(\gamma, U)$ where $\gamma \in \Gamma$ and $U$ assigns to every vertex of an independent set $S \in  \mathcal I(\mathcal G)$ a nonempty tuple $u_i$ of the above form, for each $i\in S$.
For any such pair $(\gamma, U)$, we can define its sign and weight:
\begin{itemize}
    \item \textbf{Sign:} $\mathrm{Sign}(\gamma,U)$= $(-1)^{\sum_{i \in S} |u_i|},$ where $|u_i|$ denote the length of the tuple $u_i$
    \item \textbf{Weight:} $\mathrm{wt}(\gamma,U) = x^\gamma \prod_{i \in S} \prod_{h \in u_i} x_i^{w_i(h)}$
\end{itemize}


Let $\mathrm{Init}(\gamma)$ be the set of vertices $i$ such that some reduced word of $\gamma$ begins with a syllable from $G_i$. We can then define the following sets:
\begin{equation*}
A(\gamma, U) = \{i \in \mathrm{Init}(\gamma) : \{i\} \cup S \in \mathcal{I}(G)\}
\end{equation*}
\begin{equation*}
L(\gamma, U) = A(\gamma, U) \cup S
\end{equation*}
where $S$ is the independent set corresponding to $U$.
If $L(\gamma, U) = \emptyset$, it follows that $\gamma = 1$ and $U = \emptyset$. \textit{
We repeatedly use the following key observation without mentioning further:
let \(i\in \operatorname{Init}(\gamma)\), then there is a unique syllable
\(g_i(\gamma)\in G_i\setminus\{1\}\) such that some reduced word for
\(\gamma\) begins with \(g_i(\gamma)\). Moreover, if
\(i,j\in\operatorname{Init}(\gamma)\) and \(i\ne j\), then the vertex groups
\(G_i\) and \(G_j\) commute.
}


Let $M_\Gamma$ denote the set of all $(\gamma,U)$ such that $L(\gamma,U)\neq \emptyset$.
To proceed, fix a total order on the vertex set $I$. For any nontrivial pair $(\gamma, U)$, we define the minimal element:
\begin{equation*}
i_0 = \min L(\gamma, U).
\end{equation*}

\item[Case 1: $i_0 \in A(\gamma,U)$.] 
    Since $i_0 \in \mathrm{Init}(\gamma)$, we can fix a reduced word $w$ for $\gamma$ of the form
    \begin{equation*}
        w = (g_1, g_2, \dots, g_n), \quad \text{where } g_1 \in G_{i_0} \setminus \{1\}.
    \end{equation*}
    

    We shift this first syllable out of the group element by setting $\gamma' = g_1^{-1}\gamma$ (so that $\gamma'$ has the reduced expression $(g_2, \dots, g_n)$), and update the tuple configuration $U$ to $U'$ according to the following rules:
    \begin{itemize}
        \item If $i_0 \notin S$, we add $i_0$ to the  independent set $S$ by setting $S' = S \cup \{i_0\}$ and initializing the singleton tuple $u'_{i_0} = (g_1)$. The fact that $S' \in  \mathcal I(\mathcal G)$ follows directly from the definition of $A(\gamma,U)$.
        \item If $i_0 \in S$ and $u_{i_0} = (h_1, h_2, \dots, h_r)$, we prepend $g_1$ to the existing tuple, replacing $u_{i_0}$ with $u'_{i_0} = (g_1, h_1, h_2, \dots, h_r)$.
    \end{itemize}

    \item[Case 2: $i_0 \notin A(\gamma,U)$.] 
    By the definition of $L(\gamma, U)$, this assumption implies that $i_0 \in S$. Let $u_{i_0} = (h_1, h_2, \dots, h_r)$. We shift the first element $h_1$ from the tuple to the front of the group element by setting 
    \begin{equation*}
        \gamma' = h_1\gamma.
    \end{equation*}
    The tuple configuration is then updated to $U'$ by altering the component at $i_0$:
    \begin{itemize}
        \item If $r > 1$, we replace $u_{i_0}$ with the truncated tuple $u'_{i_0} = (h_2, \dots, h_r)$.
        \item If $r = 1$, we remove the vertex $i_0$ from the active independent set, mapping $S' = S \setminus \{i_0\}$. Also this $S'\in \mathcal{I(G)}$.
    \end{itemize}
    To see that this step preserves the structural properties of our framework, let $w = (g_1, g_2, \dots, g_n)$ be our fixed reduced word for $\gamma$. We claim that the concatenated word $w' = (h_1, g_1, g_2, \dots, g_n)$ is a reduced word for $\gamma'$. Indeed, if $w'$ were not reduced, there would exist some index $k$ such that $g_k \in G_{i_0} \setminus \{1\}$ commutes with all $g_j$ for $1 \le j < k$. This would imply $i_0 \in \mathrm{Init}(\gamma)$. Because $i_0 \in S$ and $S \in  \mathcal I(\mathcal G)$, it follows that $\{i_0\} \cup S = S \in  \mathcal I(\mathcal G)$, meaning $i_0 \in A(\gamma,U)$, which  contradicts our hypothesis. Consequently, $w'$ is a valid reduced expression for $\gamma'$, and we obtain:
    \begin{equation*}
        x^{\gamma'} = x_{i_0}^{w_{i_0}(h_1)} x^\gamma.
    \end{equation*}


In both cases, let $U'$ denote the resulting tuple configuration, which is now indexed by the updated independent set $S' \in \mathcal{I}(\mathcal{G})$. We define the map $\Phi:M_\Gamma\to M_\Gamma$ by 
\begin{equation*}
    \Phi(\gamma,U) = (\gamma',U').
\end{equation*}
To see that $\Phi$ reverses the sign, observe that in each subcase, the total number of elements across all active tuples changes by exactly $\pm 1$. In Case 1, a syllable ($g_1$) is appended to the configuration $U'$, either by activating a new vertex index or by increasing the length of an existing tuple. Conversely, in Case 2, a syllable ($h_1$) is removed from the configuration $U$, either by tuple truncation or by the complete deactivation of the vertex index $i_0$. It follows that $\sum_{i \in S'} |u'_i| = \sum_{i \in S} |u_i| \pm 1$, yielding
\begin{equation*}
\operatorname{Sign}(\gamma',U') = (-1)^{\sum_{i \in S'} |u'_i|} = (-1)^{\left(\sum_{i \in S} |u_i|\right) \pm 1} = -\operatorname{Sign}(\gamma,U).
\end{equation*}

To demonstrate weight preservation, we monitor the simultaneous changes in the group element monomial and the tuple product. In Case 1, the removal of the initial syllable $g_1 \in G_{i_0}$ from $\gamma$ reduces the group element weight by a factor of $x_{i_0}^{w_{i_0}(g_1)}$, so that $x^{\gamma'} = x_{i_0}^{-w_{i_0}(g_1)} x^\gamma$. Concurrently, the addition of $g_1$ to the tuple configuration at $i_0$ scales the tuple product precisely by $x_{i_0}^{w_{i_0}(g_1)}$. Combining these modifications, we obtain
\begin{equation*}
\operatorname{wt}(\gamma',U') = \left(x_{i_0}^{-w_{i_0}(g_1)} x^\gamma\right) \cdot \left( x_{i_0}^{w_{i_0}(g_1)} \prod_{i \in S} \prod_{h \in u_i} x_i^{w_i(h)} \right) = \operatorname{wt}(\gamma,U).
\end{equation*}
In Case 2, shifting $h_1 \in G_{i_0}$ to the front of the group element scales its weight by a factor of $x_{i_0}^{w_{i_0}(h_1)}$, yielding $x^{\gamma'} = x_{i_0}^{w_{i_0}(h_1)} x^\gamma$. Concurrently, removing $h_1$ from the tuple configuration eliminates its corresponding weight from the product, which implies
\begin{equation*}
\operatorname{wt}(\gamma',U') = \left(x_{i_0}^{w_{i_0}(h_1)} x^\gamma\right) \cdot \left( x_{i_0}^{-w_{i_0}(h_1)} \prod_{i \in S} \prod_{h \in u_i} x_i^{w_i(h)} \right) = \operatorname{wt}(\gamma,U).
\end{equation*}

\smallskip
\noindent\textit{  Proof of $L(\gamma', U') = L(\gamma, U)$ under Case 1:}
Suppose $(\gamma, U)$ satisfies the conditions of Case 1, where $i_0 \in A(\gamma, U)$. We claim that $L(\gamma', U') = L(\gamma, U)$. If $i = i_0$, then $i_0 \in L(\gamma, U)$ by assumption, and because Case 1 maps $i_0$ directly into $S'$, we have $i_0 \in S' \subseteq L(\gamma', U')$. For any vertex $i \neq i_0$, we verify equality via double inclusion:
\begin{itemize}
    \item \textbf{Inclusion $L(\gamma, U) \subseteq L(\gamma', U')$:} Let $i \in L(\gamma, U) \setminus \{i_0\}$. If $i \in S$, then since $S \subseteq S'$, it follows immediately that $i \in S' \subseteq L(\gamma', U')$. If $i \in A(\gamma, U) \setminus \{i_0\}$, then $i \in \operatorname{Init}(\gamma)$ and $\{i\} \cup S \in \mathcal{I}(\mathcal{G})$. Thus we can find a reduced word of $\gamma'$ such that it begins with $g_k$ where $g_k\in G_i$, so $i \in \operatorname{Init}(\gamma')$. Because $i$ and $i_0$ are distinct elements of $\operatorname{Init}(\gamma)$, the vertices $i$ and $i_0$ must mutually commute, meaning $\{i, i_0\} \in \mathcal{I}(\mathcal{G})$.  Since $\{i\} \cup S \in \mathcal{I}(\mathcal{G})$ and $\{i, i_0\} \in \mathcal{I}(\mathcal{G})$, we have $\{i\} \cup S' = \{i\} \cup S \cup \{i_0\} \in \mathcal{I}(\mathcal{G})$. Thus, $i \in A(\gamma', U') \subseteq L(\gamma', U')$.

    \medskip
    \item \textbf{Inclusion $L(\gamma', U') \subseteq L(\gamma, U)$:} Let $i \in L(\gamma', U') \setminus \{i_0\}$. If $i \in S'$, then since $i \neq i_0$, it must be that $i \in S \subseteq L(\gamma, U)$. If $i \in A(\gamma', U') \setminus \{i_0\}$, then $i \in \operatorname{Init}(\gamma')$ and $\{i\} \cup S' \in \mathcal{I}(\mathcal{G})$. Since $i_0 \in S'$, the stability of the subset $\{i\} \cup S'$ implies that $i$ and $i_0$ commute. Hence, the syllable $g_k \in G_i$ can be moved past $g_1$ to the front of $\gamma = g_1 \gamma'$, proving $i \in \operatorname{Init}(\gamma)$. Finally, the containment $S \subseteq S'$ yields $\{i\} \cup S \in \mathcal{I}(\mathcal{G})$, which means $i \in A(\gamma, U) \subseteq L(\gamma, U)$.
\end{itemize}

\smallskip
\noindent\textit{ Proof of $L(\gamma', U') = L(\gamma, U)$  under Case 2:}
Suppose instead that $(\gamma, U)$ satisfies the conditions of Case 2, where $i_0 \in S \setminus A(\gamma, U)$. We claim that $L(\gamma', U') = L(\gamma, U)$. For the minimal index itself, $i_0 \in S \subseteq L(\gamma, U)$ by definition. Under the map, if $r > 1$, then $i_0 \in S' = S \subseteq L(\gamma', U')$. If $r = 1$, then $i_0 \notin S'$, but because the syllable $h_1 \in G_{i_0}$ is shifted to the front of the group element, $i_0 \in \operatorname{Init}(\gamma')$. Since $\{i_0\} \cup S' = S \in \mathcal{I}(\mathcal{G})$, we obtain $i_0 \in A(\gamma', U') \subseteq L(\gamma', U')$. For any vertex $i \neq i_0$, we verify the equality via double inclusion:
\begin{itemize}
    \item \textbf{Inclusion $L(\gamma, U) \subseteq L(\gamma', U')$:} Let $i \in L(\gamma, U) \setminus \{i_0\}$. If $i \in S$, then since $i \neq i_0$ and $S' \in \{S, S \setminus \{i_0\}\}$, it remains true that $i \in S' \subseteq L(\gamma', U')$. If $i \in A(\gamma, U) \setminus \{i_0\}$, then $i \in \operatorname{Init}(\gamma)$ and $\{i\} \cup S \in \mathcal{I}(\mathcal{G})$. Since $i_0 \in S$, the independence condition forces $i$ and $i_0$ to commute, implying the syllable of $i$ commutes with $h_1 \in G_{i_0}$. Thus, $i \in \operatorname{Init}(\gamma' = h_1 \gamma)$. Furthermore, because $S' \subseteq S$, the relation $\{i\} \cup S \in \mathcal{I}(\mathcal{G})$ guarantees $\{i\} \cup S' \in \mathcal{I}(\mathcal{G})$, which proves $i \in A(\gamma', U') \subseteq L(\gamma', U')$.

    \medskip
    \item \textbf{Inclusion $L(\gamma', U') \subseteq L(\gamma, U)$:} Let $i \in L(\gamma', U') \setminus \{i_0\}$. If $i \in S'$, then since $S' \subseteq S$, we immediately have $i \in S \subseteq L(\gamma, U)$. If $i \in A(\gamma', U') \setminus \{i_0\}$, then $i \in \operatorname{Init}(\gamma')$ and $\{i\} \cup S' \in \mathcal{I}(\mathcal{G})$. In Case 2, we know that $i_0$ always belongs to $ \operatorname{init}(\gamma')$. This forces $i$ to commute with $i_0$, therefore we have  $i \in \operatorname{Init}(\gamma = h_1^{-1} \gamma')$. To see that $\{i\} \cup S \in \mathcal{I}(\mathcal{G})$, note that if $r > 1$, then $S = S'$ and the condition is trivial. If $r = 1$, then $S = S' \cup \{i_0\}$; since $\{i\} \cup S' \in \mathcal{I}(\mathcal{G})$ and $i$ commutes with $i_0$, it follows that $\{i\} \cup S \in \mathcal{I}(\mathcal{G})$. Thus, $i \in A(\gamma, U) \subseteq L(\gamma, U)$.
\end{itemize}
Thus, the identity $L(\gamma', U') = L(\gamma, U)$ holds universally.

To complete the argument, it remains to verify that the map $\Phi$ is an involution on $M_\Gamma$. Suppose first that a pair $(\gamma,U) \in M_\Gamma$ satisfies the conditions of Case~1, where $i_0 = \min L(\gamma,U) \in A(\gamma,U)$. As established above, we have  $L(\gamma', U') = L(\gamma, U)$, and hence $\min L(\gamma', U') = i_0$. This immediately implies that $i_0 \notin A(\gamma',U')$. Consequently, the image $(\gamma', U')$ satisfies the precise defining hypotheses of Case~2 with respect to the same minimal active index $i_0$. Applying the rules of Case~2 to $(\gamma', U')$ shifts the first element of the tuple at $i_0$ (which is $g_1$) back to the front of the group element, setting $\gamma'' = g_1\gamma' = \gamma$, and perfectly reverses the structural modification of the tuple configuration. This yields $\Phi(\gamma', U') = (\gamma, U)$.

Conversely, suppose that $(\gamma,U)$ satisfies the conditions of Case~2, so that the minimal index satisfies $i_0 \in S \setminus A(\gamma,U)$. In the resulting pair $(\gamma',U')$, the minimal index is again preserved. Moreover, because the initial element $h_1 \in G_{i_0}$ of the tuple is shifted to the front of the group element, we have $i_0 \in \operatorname{Init}(\gamma')$. Since $\{i_0\} \cup S' = S \in \mathcal{I}(\mathcal{G})$, it follows that $i_0 \in A(\gamma',U')$. Therefore, the pair $(\gamma',U')$ satisfies the hypotheses of Case~1. Applying the rules of Case~1 to $(\gamma',U')$ removes this initial syllable $h_1$ from $\gamma'$ and prepends it back to the tuple configuration at $i_0$, which exactly recovers the original pair; that is, $\Phi(\gamma', U') = (\gamma, U)$.

Consequently, $\Phi$ defines a sign-reversing, weight-preserving involution on the set of all pairs $(\gamma, U) \neq (1, \emptyset)$. It follows that all nontrivial terms cancel out in pairs within the summation. The only surviving contribution arises from the pair $(1, \emptyset)$, whose weight is exactly $1$. This completes the proof.
\end{proof}

\end{document}
